\section{Results}
\label{sec:results}
\label{sec:experiments}

In the following, we introduce each benchmark, explain its relevance to
particle-based sampling, and report the corresponding empirical results. The
benchmark suite consists of a 50D Gaussian target, 2D synthetic targets, a
multimodal Gaussian mixture (Mix8), UCI logistic regression, a Bayesian neural
network, and large-scale hierarchical logistic regression.

\subsection{50D Gaussian}

This benchmark tests the stability of sampling in a high-dimensional setting with anisotropic
covariance structure. We run 1,000 sampling iterations with checkpointing every 50 steps
and a 5-seed ablation. Reported metrics are moment errors
($\|\hat{\mu}-\mu\|_2$ and $\|\hat{\Sigma}-\Sigma\|_F$), ESS, KSD, and time-to best-checkpoint.
\begin{figure}[t]
  \centering
  \safeincludegraphics[width=0.9\linewidth]{../figures/50d/summary_metric_panels.png}
  \caption{50D Gaussian final-checkpoint summary across methods. The four panels
  report $\|\hat{\mu}-\mu\|_2$, $\|\hat{\Sigma}-\Sigma\|_F$, KSD, and ESS,
  respectively; markers show mean values across seeds and error bars indicate
  one standard deviation. Lower is better for the first three metrics, whereas
  higher is better for ESS. This figure highlights robustness differences under
  strong anisotropy: Adapt-MR-SVGD is the only particle variant that keeps both
  moment errors and KSD under control.}
  \label{fig:gauss50_summary}
\end{figure}

\begin{figure}[t]
  \centering
  \safeincludegraphics[width=0.9\linewidth]{../figures/50d/pareto_mu_cov_final_mean_err.png}
  \caption{50D Gaussian: final mean $\pm$ std Pareto plot in moment-error space
  with marker size encoding ESS (left) and wall time (right).}
  \label{fig:gauss50_pareto}
\end{figure}

Among the methods, \cref{fig:gauss50_summary} shows that
Adapt-MR-SVGD is the only variant that keeps both moment errors and KSD under
control in this anisotropic Gaussian setting: it substantially improves on
MR-SVGD and avoids the severe degradation exhibited by vanilla SVGD and
Strang-SVGD. The chain baselines remain competitive on some final-checkpoint
fidelity metrics, but they do so with ESS near $4$, far below the particle
methods. Within the particle family, \cref{fig:gauss50_pareto} shows that MR-SVGD offers
a cheaper but less accurate baseline, whereas Adapt-MR-SVGD
achieves the strongest quality at a higher computational cost.

\subsection{2D synthetic targets}
In this benchmark, we evaluate a suite of 2D targets with varying geometry and multimodality, including banana, ring, squiggly, two-moons, and funnel.
To contextualize the geometry of these targets, \cref{fig:2d_demo_panel} shows
short-run visualization panels for representative cases.
The metrics are averaged over five
seeds, each with 1,000 iterations and checkpointing every 20 steps. Here KSD serves
both as the stop monitor and the primary quality metric. We add a safety guard
against non-finite KSD values in the stopping policy, because several targets
can produce unstable particle trajectories. We report mean log probability as a secondary
density-concentration score, gradient/kernel evaluations and wall time as cost
metrics, and ESS as a supplemental mixing proxy. Moment errors and grid-based
$L_1$ discrepancy are retained as diagnostics, but they are not used in the
main 2D summary table in \cref{tab:2d_summary}.

\begin{figure}[t]
  \centering
  \begin{subfigure}[t]{0.95\linewidth}
    \centering
    \safeincludegraphics[width=\linewidth]{../figures/2d_demo_tuned/banana_demo.png}
    \caption{Banana target: initial and short-run final particle states.}
    \label{fig:2d_demo_banana}
  \end{subfigure}

  \vspace{0.35em}

  \begin{subfigure}[t]{0.95\linewidth}
    \centering
    \safeincludegraphics[width=\linewidth]{../figures/2d_demo_tuned/squiggly_demo.png}
    \caption{Squiggly target: initial and short-run final particle states.}
    \label{fig:2d_demo_squiggly}
  \end{subfigure}
  \caption{2D target visualization panels produced with Adapt-MR-SVGD under
  visualization-only settings. Each subfigure shows target-density contours with
  initial particles (left) and short-run final particles (right).}
  \label{fig:2d_demo_panel}
\end{figure}

\begin{figure}[t]
  \ContinuedFloat
  \centering
  \begin{subfigure}[t]{0.95\linewidth}
    \centering
    \safeincludegraphics[width=\linewidth]{../figures/2d_demo_tuned/two_moons_demo.png}
    \caption{Two-moons target: initial and short-run final particle states.}
    \label{fig:2d_demo_two_moons}
  \end{subfigure}

  \vspace{0.35em}

  \begin{subfigure}[t]{0.95\linewidth}
    \centering
    \safeincludegraphics[width=\linewidth]{../figures/2d_demo_tuned/ring_demo.png}
    \caption{Ring target: initial and short-run final particle states.}
    \label{fig:2d_demo_ring}
  \end{subfigure}
  \caption[]{2D target visualization panels (continued).}
\end{figure}

\begin{table}[t]
  \centering
  \scriptsize
  \setlength{\tabcolsep}{3pt}
  \begin{tabular}{lcccccc}
    \toprule
    Method & Fails$\downarrow$ & KSD$\downarrow$ & mean logp$\uparrow$ & ESS$\uparrow$ & grad evals$\downarrow$ & wall (s)$\downarrow$ \\
    \midrule
    Adapt-MR-SVGD & 0 & 0.308 & -0.175 & 128.0 & 182 & 0.79 \\
    MR-SVGD & 0 & 40.13 & -278.0 & 121.8 & 80 & 1.13 \\
    SVGD & 4 & 26.02 & $-1.93\times 10^4$ & 128.0 & 80 & 0.32 \\
    Strang-SVGD & 4 & 30.09 & $-1.20\times 10^5$ & 121.7 & 121 & 1.55 \\
    \bottomrule
  \end{tabular}
  \caption{2D targets quality-cost summary across particle methods. Entries are
  medians over the $25$ target-seed runs, using the best checkpoint selected by
  the KSD monitor for each run. ``Fails'' counts runs whose KSD monitor never
  attained a finite value; methods are ordered first by failure count and then
  by median KSD. KSD is the primary distributional quality metric; mean log
  probability is reported as a secondary score, ESS as a supplemental mixing
  indicator, and gradient evaluations / wall time as cost metrics. Moment
  errors and grid $L_1$ remain diagnostic quantities and are omitted from this
  main summary. The main pattern is that Adapt-MR-SVGD is the only zero-failure
  particle method and gives the best overall KSD across the suite, while banana
  shows little separation and funnel remains a nuanced exception.}
  \label{tab:2d_summary}
\end{table}

Across the full 2D suite, \cref{tab:2d_summary} shows that Adapt-MR-SVGD gives
the strongest overall distributional fidelity: it is the only particle method
with zero failures, attains the best overall median KSD, and also attains the
best median mean log-density while retaining high ESS. The gain is
concentrated on the geometry-driven ring, squiggly, and two-moons targets,
where Adapt-MR-SVGD
achieves the best mean KSD over seeds (approximately $0.231$, $1.03$, and
$0.308$, respectively) and the single-rate particle methods degrade sharply.
On the easier banana target, by contrast, all particle methods attain nearly
identical KSD values around $0.116$, consistent with the comparatively mild
geometry illustrated in \cref{fig:2d_demo_panel}; here multirate treatment is
not especially beneficial.

The funnel target remains the main exception. Fixed MR-SVGD attains the best
mean KSD on funnel ($0.261$ versus $0.456$ for Adapt-MR-SVGD), but
Adapt-MR-SVGD still remains fully stable and yields substantially better mean
log-density ($-0.023$ versus $-2.89$). Vanilla SVGD and Strang-SVGD each fail
on four of the five funnel seeds, so finite-run medians alone would understate
their instability; this is why \cref{tab:2d_summary} reports failure counts
alongside the quality and cost summaries. Overall, these 2D results sharpen the
main qualitative claim: adaptive multirate integration is most valuable when
the interaction between drift and repulsion becomes strongly geometry-dependent,
rather than as a uniform improvement on every benign target. In this setting,
ESS is much less discriminative than KSD and mean log-density, so we treat it
as a supplemental mixing indicator rather than a primary quality metric.

\subsection{Mixture2D (mix8)}

This benchmark emphasizes multimodal exploration and mass allocation across
modes. All methods start from the same small Gaussian centered at the origin,
away from the modal ring, and then run for a fixed budget of 1,000 iterations
with checkpointing every 20 steps across 5 seeds. For this benchmark we use a
shared multiscale RBF kernel with bandwidth factors $\{0.5,1,2\}$ for all
particle methods so that local spacing and longer-range repulsion are both
represented on the modal ring. We treat this as an experimental kernel-choice
ablation rather than a change to the multirate update itself; prior SVGD work
has also considered broader kernel families beyond a single scalar RBF
\cite{wang2019matrixsvgd,stordal2021pkernel}. Let $p_k$ denote the empirical
mass assigned to mixture center $k$ by nearest-center assignment. We define
mode coverage as the fraction of modes with $p_k\ge 0.05$, normalized mode
entropy as $-\sum_{k=1}^{K} p_k\log p_k / \log K$, and mode imbalance as the
standard deviation of $(p_1,\ldots,p_K)$. The primary metrics are therefore
mode coverage, normalized mode entropy, and mode imbalance; KSD and mean log
probability are retained as secondary diagnostics, with ESS reported as a
supplemental mixing proxy.

\begin{figure}[t]
  \centering
  \safeincludegraphics[width=0.9\linewidth]{../figures/mixture2d/mix8/mix8_summary_metric_panels.png}
  \caption{Mixture2D (mix8) fixed-budget final-checkpoint summary across methods.
  The four
  panels report mode coverage, mode entropy, mode imbalance, and KSD;
  markers show mean values across seeds and error bars indicate one standard
  deviation. Higher is better for coverage and entropy, while lower is better
  for mode imbalance and KSD.}
  \label{fig:mix8_summary}
\end{figure}

\begin{table}[t]
  \centering
  \small
  \setlength{\tabcolsep}{7pt}
  \caption{Mixture2D (mix8) mean final metrics over five seeds at the fixed
  1,000-iteration budget. Higher is better for coverage and entropy; lower is
  better for KSD. We omit minimum mass because it is zero for all methods at
  this budget. This table summarizes the fixed-budget multimodal result:
  Adapt-MR-SVGD gives the strongest overall exploration, while vanilla and
  Strang spread more than MR-SVGD but remain far less stable by KSD.}
  \begin{tabular}{lrrr}
    \toprule
    Method & Coverage$\uparrow$ & Entropy$\uparrow$ & KSD$\downarrow$ \\
    \midrule
    Adapt-MR-SVGD & 0.500 & 0.544 & 0.360 \\
    SVGD          & 0.375 & 0.453 & 3162.9 \\
    Strang-SVGD   & 0.325 & 0.403 & 3883.0 \\
    MR-SVGD       & 0.250 & 0.238 & 74.3 \\
    SGHMC         & 0.125 & 0.000 & 2.29 \\
    SGLD          & 0.125 & 0.000 & 2.53 \\
    \bottomrule
  \end{tabular}
  \label{tab:mix8_summary}
\end{table}

Under this harder off-mode initialization, adaptive MR-SVGD still gives the
strongest overall multimodal exploration at the fixed budget, now reaching the
highest mean final mode coverage and mode entropy across seeds
($0.500$ and $0.544$, respectively); \cref{tab:mix8_summary} provides the
corresponding mean-over-seeds summary, while \cref{fig:mix8_summary} carries
the cross-seed variability. The shared multiscale kernel also makes the
single-rate particle baselines more competitive than under the plain single-RBF
setting: vanilla SVGD and Strang-SVGD both cover more modes and attain higher
mode entropy, but their KSD values remain far worse than the adaptive method,
so the extra spread still comes with much weaker stability. Fixed MR-SVGD
remains stable and slightly improves its KSD over the single-RBF version, but
it does not match the adaptive variant on the main multimodal metrics. The
imbalance panel in \cref{fig:mix8_summary} complements coverage and entropy by
showing that the adaptive method also allocates mass more evenly across the
discovered modes. SGLD and SGHMC remain trapped in a single mode, with final
coverage fixed at $0.125$ and zero mode entropy. The minimum mass statistic is
still zero for all methods at this budget, so in this redesigned benchmark
coverage, entropy, and imbalance are the informative multimodal diagnostics
while KSD remains a secondary stability measure.

\subsection{UCI logistic regression}

We benchmark Bayesian logistic regression on breast\_cancer, ionosphere,
spambase, and a5a, using a Gaussian prior on weights (as implemented in the
benchmark script) and 5 seeds per dataset. We report test accuracy, NLL, ECE,
ESS, and mean log probability; KSD is optional and disabled by default in this
benchmark configuration. Results are summarized with per-dataset panels in
\cref{fig:uci_summary} and a grouped best-NLL summary table (particle methods
only) in \cref{tab:uci_nll_summary}; both use the best finite-NLL checkpoint
selected under the global predictive stopping policy
\cite{dua2019uci}.

\begin{figure}[t]
  \centering
  \begin{subfigure}[t]{0.9\linewidth}
    \centering
    \safeincludegraphics[width=\linewidth]{../figures/uci/breast_cancer/breast_cancer_summary_metric_panels.png}
    \caption{breast\_cancer}
    \label{fig:uci_breast_cancer}
  \end{subfigure}

  \vspace{0.25em}

  \begin{subfigure}[t]{0.9\linewidth}
    \centering
    \safeincludegraphics[width=\linewidth]{../figures/uci/ionosphere/ionosphere_summary_metric_panels.png}
    \caption{ionosphere}
    \label{fig:uci_ionosphere}
  \end{subfigure}
  \caption{UCI logistic regression summary across datasets. Each subpanel
  reports the best-checkpoint mean across seeds with one standard deviation for
  test accuracy, NLL, ECE, and ESS, where the checkpoint is selected by the
  NLL-based early-stop monitor. Higher is better for accuracy and ESS; lower is
  better for NLL and ECE. Remaining datasets are shown in the continued panel.}
  \label{fig:uci_summary}
\end{figure}

\begin{figure}[t]
  \ContinuedFloat
  \centering
  \begin{subfigure}[t]{0.9\linewidth}
    \centering
    \safeincludegraphics[width=\linewidth]{../figures/uci/spambase/spambase_summary_metric_panels.png}
    \caption{spambase}
    \label{fig:uci_spambase}
  \end{subfigure}

  \vspace{0.25em}

  \begin{subfigure}[t]{0.9\linewidth}
    \centering
    \safeincludegraphics[width=\linewidth]{../figures/uci/a5a/a5a_summary_metric_panels.png}
    \caption{a5a}
    \label{fig:uci_a5a}
  \end{subfigure}
  \caption[]{UCI logistic regression summary across datasets (continued).}
\end{figure}

\begin{table}[t]
  \centering
  \scriptsize
  \setlength{\tabcolsep}{4pt}
  \begin{tabular}{lccc}
    \toprule
    Method & Best NLL$\downarrow$ & Acc@best NLL$\uparrow$ & Wall@best (s)$\downarrow$ \\
    \midrule
    \multicolumn{4}{l}{\textit{breast\_cancer}} \\
    Adapt-MR-SVGD & 0.062 & 0.986 & 8.9 \\
    MR-SVGD & 0.089 & 0.981 & 6.2 \\
    Strang-SVGD & 0.094 & 0.979 & 6.5 \\
    SVGD & 0.104 & 0.979 & 2.0 \\
    \midrule
    \multicolumn{4}{l}{\textit{ionosphere}} \\
    Adapt-MR-SVGD & 0.383 & 0.900 & 3.4 \\
    SVGD & 0.438 & 0.906 & 0.9 \\
    Strang-SVGD & 0.447 & 0.897 & 5.2 \\
    MR-SVGD & 0.500 & 0.897 & 6.1 \\
    \midrule
    \multicolumn{4}{l}{\textit{spambase}} \\
    SVGD & 0.288 & 0.918 & 1.3 \\
    Strang-SVGD & 0.288 & 0.924 & 8.8 \\
    Adapt-MR-SVGD & 0.361 & 0.882 & 8.1 \\
    MR-SVGD & 0.467 & 0.877 & 4.3 \\
    \midrule
    \multicolumn{4}{l}{\textit{a5a}} \\
    Adapt-MR-SVGD & 0.453 & 0.778 & 36.5 \\
    Strang-SVGD & 0.471 & 0.836 & 3.0 \\
    MR-SVGD & 0.474 & 0.809 & 14.6 \\
    SVGD & 0.478 & 0.835 & 2.6 \\
    \bottomrule
  \end{tabular}
  \caption{UCI logistic-regression particle-method summary over five seeds. For
  each dataset and method, we select the best finite-NLL checkpoint per seed
  and report the method-wise mean of that best NLL, the accuracy at that
  checkpoint, and the corresponding wall time. ECE is summarized in the
  per-dataset panels of \cref{fig:uci_summary}.}
  \label{tab:uci_nll_summary}
\end{table}

UCI logistic regression serves mainly as a predictive sanity check rather than
as the flagship evidence for multirate gains. The rankings are mixed.
Adapt-MR-SVGD gives the best overall NLL on breast\_cancer, while on
ionosphere, spambase, and a5a either SGHMC or a single-rate particle method is
competitive or superior. This weaker separation is consistent with the main
claim of the paper: multirate benefits are most pronounced on anisotropic,
multimodal, or hierarchical posteriors, whereas lower-dimensional logistic
regression tasks show smaller and less uniform gains. On the smaller datasets,
accuracy is also tightly quantized, so NLL is the more discriminative summary.

\subsection{Bayesian neural network}

We evaluate a one-hidden-layer BNN (hidden width 32) on spambase and a5a,
using a Gaussian prior, a fixed split seed, and 5 sampler seeds per dataset.
For each dataset, we carve a validation split from the training data, select
checkpoints by validation NLL under the global predictive stopping policy, and
report test accuracy, test NLL, test ECE, ESS, and mean log probability at the
selected checkpoint. KSD is optional and disabled by default in this
configuration.

\begin{figure}[t]
  \centering
  \begin{subfigure}[t]{0.9\linewidth}
    \centering
    \safeincludegraphics[width=\linewidth]{../figures/bnn/spambase/spambase_summary_metric_panels.png}
    \caption{spambase}
    \label{fig:bnn_spambase}
  \end{subfigure}

  \vspace{0.25em}

  \begin{subfigure}[t]{0.9\linewidth}
    \centering
    \safeincludegraphics[width=\linewidth]{../figures/bnn/a5a/a5a_summary_metric_panels.png}
    \caption{a5a}
    \label{fig:bnn_a5a}
  \end{subfigure}
  \caption{BNN summary across datasets. Each subpanel reports the
  best-checkpoint mean across seeds with one standard deviation for test
  accuracy, NLL, ECE, and ESS, where the checkpoint is selected by the
  validation-NLL early-stop monitor. Higher is better for accuracy and ESS;
  lower is better for NLL and ECE.}
  \label{fig:bnn_summary}
\end{figure}

\begin{table}[t]
  \centering
  \scriptsize
  \setlength{\tabcolsep}{4pt}
  \begin{tabular}{lccc}
    \toprule
    Method & Best NLL$\downarrow$ & Acc@best NLL$\uparrow$ & Wall@best (s)$\downarrow$ \\
    \midrule
    \multicolumn{4}{l}{\textit{spambase}} \\
    Adapt-MR-SVGD & 0.174 & 0.940 & 35.7 \\
    Strang-SVGD & 0.187 & 0.930 & 51.2 \\
    SVGD & 0.187 & 0.930 & 21.1 \\
    MR-SVGD & 0.188 & 0.931 & 51.5 \\
    \midrule
    \multicolumn{4}{l}{\textit{a5a}} \\
    MR-SVGD & 0.334 & 0.844 & 5.3 \\
    Adapt-MR-SVGD & 0.347 & 0.836 & 8.1 \\
    SVGD & 0.362 & 0.834 & 17.5 \\
    Strang-SVGD & 0.378 & 0.832 & 49.6 \\
    \bottomrule
  \end{tabular}
  \caption{BNN particle-method summary over five seeds. For each dataset and
  method, we select the best finite-validation-NLL checkpoint per seed and
  report the method-wise mean of the corresponding test NLL, test accuracy, and
  wall time. ECE is summarized in the per-dataset panels of
  \cref{fig:bnn_summary}.}
  \label{tab:bnn_nll_summary}
\end{table}

The two BNN datasets give a more informative but less uniform predictive story
than the earlier small-data configuration. On spambase, Adapt-MR-SVGD attains
the best overall test NLL ($0.174$) and the highest mean test accuracy
($0.940$), with the remaining particle methods clustered together around
$0.187$. On a5a, the ranking shifts: MR-SVGD attains the best overall test NLL
($0.334$) and the highest mean test accuracy ($0.844$), Adapt-MR-SVGD remains
the second-best particle method ($0.347$), and the single-rate particle
variants are again worse. The chain baselines are competitive in isolated
cases---SGHMC on spambase and SGLD on a5a---but the multirate particle methods
form the strongest pair overall across the two datasets. This makes the BNN
benchmark corroborative rather than primary evidence: multirate structure improves
predictive quality over the single-rate particle baselines, but the preferred
multirate variant depends on the dataset.

\subsection{Large-scale hierarchical logistic regression}

To stress-test multirate behavior in a high-dimensional anisotropic posterior,
we use a large-scale hierarchical logistic regression benchmark with grouped
random effects. The model combines high parameter dimension with funnel-like
geometry induced by a global scale parameter, making it a natural benchmark for
joint quality-cost assessment of SVGD variants.

\paragraph{Problem definition}
Given binary outcomes $y_i\in\{0,1\}$, sparse features $x_i\in\mathbb{R}^p$, and
group indices $g_i\in\{1,\ldots,G\}$, we define
\begin{equation}
\begin{aligned}
y_i \mid x_i,g_i,\beta,\alpha,u
&\sim \operatorname{Bernoulli}\!\left(
\sigma\!\left(\alpha + x_i^\top\beta + u_{g_i}\right)\right), \\
&\qquad i=1,\ldots,n,
\end{aligned}
\label{eq:hlr-likelihood}
\end{equation}
where $\sigma(\cdot)$ is the logistic link. We use a non-centered
parameterization for group effects,
\begin{equation}
\begin{aligned}
\beta &\sim \mathcal{N}(0,\sigma_\beta^2 I_p), &
\alpha &\sim \mathcal{N}(0,\sigma_\alpha^2), \\
u_j &= \tau z_j, &
z_j &\sim \mathcal{N}(0,1), \\
\log\tau &\sim \mathcal{N}(\mu_\tau,\sigma_\tau^2),
\end{aligned}
\label{eq:hlr-priors}
\end{equation}
for $j=1,\ldots,G$. The posterior (up to normalization) is
\begin{equation}
\begin{aligned}
\pi(\theta) &\propto
\left[\prod_{i=1}^{n}
\operatorname{Bernoulli}\!\left(
y_i;\sigma\!\left(\alpha+x_i^\top\beta+u_{g_i}\right)\right)\right] \\
&\qquad \times p(\beta)\,p(\alpha)\,p(z)\,p(\log\tau), \\
\theta &= (\beta,\alpha,z,\log\tau).
\end{aligned}
\label{eq:hlr-posterior}
\end{equation}

\paragraph{Protocol}
For the long-tail grouped setting, we run $n=10^6$ observations with $p=300$
features and $G=50{,}000$ groups. Each method runs up to 1,000 iterations with
checkpointing every 20 iterations and 5 random seeds; particle methods use 32
particles, and single-chain baselines use a rolling-window evaluation set. This
benchmark follows the global predictive NLL-based stopping policy and adds a
two-checkpoint non-finite fail-fast guard because unstable methods can quickly
leave the finite-NLL regime. ECE is reported as a metric but is not used as a
stopping trigger.

\begin{table}[t]
  \centering
  \scriptsize
  \setlength{\tabcolsep}{4pt}
  \begin{tabular}{lcccc}
    \toprule
    Method & Finite-best$\uparrow$ & Best NLL$\downarrow$ & Acc@best NLL$\uparrow$ & Wall@best (s)$\downarrow$ \\
    \midrule
    Adapt-MR-SVGD & 5/5 & 0.613 & 0.658 & 123.3 \\
    MR-SVGD & 5/5 & 0.623 & 0.656 & 92.5 \\
    SGHMC & 2/5 & 0.860 & 0.560 & 77.9 \\
    SVGD & 3/5 & 0.885 & 0.577 & 154.2 \\
    Strang-SVGD & 3/5 & 1.263 & 0.579 & 163.6 \\
    SGLD & 4/5 & 1.823 & 0.576 & 101.6 \\
    \bottomrule
  \end{tabular}
  \caption{HLR long-tail summary over five seeds. For each seed and method, we
  select the best finite-NLL checkpoint; the table reports method-wise means of
  that best NLL, the accuracy at that checkpoint, and the corresponding wall
  time. ``Seeds with finite best''
  counts runs that reached at least one finite-NLL checkpoint before stopping.
  This table is the main hierarchical stress test: Adapt-MR-SVGD achieves the
  best best-NLL while remaining finite on all seeds, and MR-SVGD is the
  closest lower-cost baseline.}
  \label{tab:hlr_longtail_summary}
\end{table}

\paragraph{Result summary}
\cref{tab:hlr_longtail_summary} shows that Adapt-MR-SVGD achieves the best mean
best-NLL while remaining finite across all seeds. MR-SVGD is a close second in
NLL and reaches its best checkpoint faster on average (92.5\,s versus
123.3\,s). SVGD/Strang-SVGD and SGHMC show frequent non-finite behavior in this
regime, while SGLD remains finite more often but with substantially worse NLL.

\begin{table}[t]
  \centering
  \scriptsize
  \setlength{\tabcolsep}{4pt}
  \begin{tabular}{lcccc}
    \toprule
    Method & Finite-best$\uparrow$ & Best NLL$\downarrow$ & Acc@best NLL$\uparrow$ & Wall@best (s)$\downarrow$ \\
    \midrule
    Adapt-MR-SVGD & 5/5 & 0.681 & 0.601 & 94.6 \\
    SVGD & 1/5 & 0.700 & 0.548 & 54.9 \\
    MR-SVGD & 5/5 & 0.706 & 0.597 & 44.2 \\
    Strang-SVGD & 1/5 & 0.708 & 0.584 & 241.7 \\
    SGLD & 5/5 & 0.743 & 0.528 & 66.0 \\
    SGHMC & 5/5 & 0.791 & 0.554 & 79.7 \\
    \bottomrule
  \end{tabular}
  \caption{HLR uniform-group summary over five seeds. As in
  \cref{tab:hlr_longtail_summary}, each entry uses the best finite-NLL
  checkpoint per seed and reports method-wise means for NLL, accuracy at that
  checkpoint, and wall time.}
  \label{tab:hlr_uniform_summary}
\end{table}

\paragraph{Uniform-group sensitivity check}
\cref{tab:hlr_uniform_summary} provides a compact robustness check under
uniform group assignments. Adapt-MR-SVGD remains best in mean best-NLL with
full finite-seed coverage (5/5), and MR-SVGD remains close while being fastest
among methods with full finite coverage. This supports the same ordering seen
in the long-tail regime while keeping the long-tail setting as the primary
stress test for our main claim.
